\chapter{Cómo usar este manual}\label{cap:introduccion}

Qadra es la mesa de trabajo del despacho: expedientes, documentos, actividad
procesal, plazos y contratos con el cliente. Este manual explica, paso a paso,
cómo realizar cada tarea frente a la pantalla. Busque su rol en la
\cref{tab:roles-manual} y lea solo las partes que le corresponden.

\begin{table}[H]
  \centering
  \caption{Qué parte del manual lee cada rol.}\label{tab:roles-manual}
  \begin{tabularx}{\textwidth}{L{3.3cm} Y L{3.6cm}}
    \toprule
    \textbf{Rol} & \textbf{Qué hace en Qadra} & \textbf{Lea} \\
    \midrule
    Administrador (titular del despacho)
      & Consulta todos los expedientes. Administra el equipo, las
        asignaciones, la bitácora y los calendarios judiciales.
      & \Cref{parte:despacho} completa. \\
    Litigante
      & Trabaja en los expedientes que tiene asignados: documentos,
        actividad procesal, contratos e informes.
      & \Cref{parte:despacho}, excepto las funciones del Administrador. \\
    Asistente legal
      & Carga y consulta documentos de sus expedientes asignados.
        Consulta la actividad procesal. No sella documentos.
      & \Cref{parte:despacho}, excepto las funciones del Administrador. \\
    Cliente
      & Vincula su firma digital, firma su contrato y valida paquetes
        de evidencia.
      & Parte B. \\
    Verificador externo (perito, juez u otra persona)
      & Comprueba un paquete de evidencia sin tener cuenta en Qadra.
      & Parte B, verificación independiente. \\
    \bottomrule
  \end{tabularx}
\end{table}

\section*{Convenciones}
\begin{itemize}
  \item Los botones, menús y campos aparecen en \textbf{negritas}, tal como
        se ven en pantalla.
  \item Los mensajes del sistema aparecen entre comillas:
        \enquote{Correo o contraseña incorrectos.}
  \item Los nombres de archivo aparecen en \texttt{tipografía fija}.
  \item Cada pantalla indica al inicio qué roles pueden usarla
        (\textbf{Disponible para:}).
\end{itemize}

\importante{Qadra es un prototipo académico. El sellado de documentos es una
demostración técnica: no equivale a una constancia NOM-151 emitida por un
Prestador de Servicios de Certificación autorizado. Los certificados provienen
de una autoridad certificadora interna de demostración y no sustituyen a la
e.firma del SAT.}
